

 title\+: Temperature Control in a Flange permalink\+: doc\+\_\+\+Temp\+Control\+Flange.\+html keywords\+: Open\+F\+O\+AM sidebar\+: doc\+\_\+sidebar \subsection*{folder\+: doc }

\subsection*{Temperature Control in a Flange}

This example demonstrates the ability to pass boundary condition data using a pregenerated F\+MU into the Open\+Foam environment using the F\+M\+U4\+F\+O\+AM extension. The F\+MU was generated using Open\+Modelica and is very similar to Modelica Standard Library\textquotesingle{}s \char`\"{}\+Controlled\+Temperature\char`\"{} example from the Heat\+Transfer library shown in the following figure.

\{\% include image.\+html file=\char`\"{}\+Temp\+Control\+Flange/original\+\_\+schematic.\+jpg\char`\"{} alt=\char`\"{}\+Schematic of the Controlled\+Temperature example in Open\+Modelica\char`\"{} caption=\char`\"{}\+Schematic of the Controlled\+Temperature example in Open\+Modelica\char`\"{} \%\}

In this F\+M\+U4\+F\+O\+AM example case, the F\+MU is used to control the temperature of a flange like geometry by imposing a heat flux boundary condition to the Open\+Foam domain in order to heat up the flange.

The schematic of the F\+MU is depicted in the following sketch\+:

\{\% include image.\+html file=\char`\"{}\+Temp\+Control\+Flange/\+Controlled\+Temperature\+Coupled.\+png\char`\"{} alt=\char`\"{}\+Schematic of the F\+M\+U in Open\+Modelica\char`\"{} caption=\char`\"{}\+Schematic of the F\+M\+U in Open\+Modelica\char`\"{} \%\}

There are two input variables to the F\+MU\+:


\begin{DoxyEnumerate}
\item The actual temperature of the patch, the boundary condition is assigned to
\item the time derivative of the temperature at that specific patch
\end{DoxyEnumerate}

The temperature is then compared to a target temperature at the current time step by an on/off controller element. If the temperature is lower than the target temperature by a predefined threshold, the temperature controller triggers a switch thus activating a heater element. The output signal of that heater element is a heat flux signal that is then used as the output signal of the F\+MU dictating the heat flux at the boundary in the Open\+Foam domain.

In this particular case, the target temperature is a rising linear function. The temperature controller is responible for applying the necessary heat flux to heat up the flange. The heat flux output is shown in the following figure.

\{\% include image.\+html file=\char`\"{}\+Temp\+Control\+Flange/\+Q.\+png\char`\"{} alt=\char`\"{}\+Output heat flux Q\char`\"{} caption=\char`\"{}\+Output heat flux Q\char`\"{} \%\}

By applying the heat flux from the previous plot to the patch of the Open\+Foam domain, the temperature of the flange follows the target function in a saw tooth pattern as expected using an on/off controller with build-\/in hysteresis.

\{\% include image.\+html file=\char`\"{}\+Temp\+Control\+Flange/\+T.\+png\char`\"{} alt=\char`\"{}\+Input and target temperature\char`\"{} caption=\char`\"{}\+Input and target temperature\char`\"{} \%\}

An animation of the surface temperature is shown here\+:

\{\% include image.\+html file=\char`\"{}\+Temp\+Control\+Flange/\+Temperature\+Controlled\+Flange.\+gif\char`\"{} alt=\char`\"{}\+Input and target temperature\char`\"{} caption=\char`\"{}\+Input and target temperature\char`\"{} \%\} 